\section{Results}
\label{sec:results}

We organize results around the four sub-hypotheses: performance comparison (H1, \secref{sec:perf}), mutual information analysis (H2--H3, \secref{sec:mi}), and ablation of temporal context features (H4, \secref{sec:ablation}).

\subsection{Performance Comparison (H1)}
\label{sec:perf}

\Tabref{tab:performance} presents the mean team reward across conditions. Temporal heterogeneity produces strikingly different effects depending on task structure.

\begin{table}[t]
    \centering
    \caption{Mean team reward ($\pm$ std) over the last 50 episodes, averaged across 5 seeds. Best result per environment in \textbf{bold}. Heterogeneity helps in complementary-role tasks (\foraging) but hurts in synchronization-dependent tasks (\relay, \rendezvous).}
    \begin{tabular}{@{}lccccc@{}}
        \toprule
        \textbf{Environment} & \textbf{\hetero} & \textbf{\homo} & \textbf{\ablation} & \textbf{$p$-value} & \textbf{Cohen's $d$} \\
        \midrule
        \relay & $1.25 \pm 1.5$ & $\mathbf{3.56 \pm 0.5}$ & $1.70 \pm 0.5$ & $0.004$ & $-2.94$ \\
        \foraging & $\mathbf{20.67 \pm 1.2}$ & $14.30 \pm 0.7$ & $20.43 \pm 1.1$ & $< 0.0001$ & $+6.56$ \\
        \rendezvous & $149.95 \pm 8.0$ & $\mathbf{180.03 \pm 2.5}$ & $139.28 \pm 6.0$ & $0.001$ & $-4.83$ \\
        \bottomrule
    \end{tabular}
    \label{tab:performance}
\end{table}

In \foraging, heterogeneous agents outperform homogeneous agents by 45\% ($p < 0.0001$, $d = 6.56$). This advantage arises from the reward structure: slow agents collect $2\times$ reward per resource, so temporal diversity naturally complements the task. In contrast, \relay shows a 65\% homogeneous advantage ($p = 0.004$, $d = -2.94$) and \rendezvous shows a 17\% advantage ($p = 0.001$, $d = -4.83$). Both tasks require precise synchronization that heterogeneous action frequencies disrupt. \textbf{H1 is partially supported}: heterogeneity helps when roles are naturally complementary but hurts when synchronization is required.

\subsection{Mutual Information Analysis (H2, H3)}
\label{sec:mi}

\para{Cross-agent MI at lag 0.} \Tabref{tab:mi} shows the instantaneous mutual information between agent value functions. Contrary to the hypothesis, homogeneous agents show \emph{higher} cross-agent MI than heterogeneous agents across all three environments. Agents operating at the same frequency naturally develop more correlated value functions because they observe similar temporal patterns. \textbf{H2 is not supported.}

\begin{table}[t]
    \centering
    \caption{Cross-agent mutual information at lag 0 (instantaneous coupling) and MI decay rate (ratio of MI at lag 5 to MI at lag 0). Higher decay rate indicates more persistent temporal dependencies. The \ablation condition shows dramatically lower MI, confirming that temporal features drive cross-agent coupling.}
    \begin{tabular}{@{}lcccccc@{}}
        \toprule
        & \multicolumn{3}{c}{\textbf{MI at lag 0}} & \multicolumn{3}{c}{\textbf{MI decay rate}} \\
        \cmidrule(lr){2-4} \cmidrule(lr){5-7}
        \textbf{Environment} & \textbf{\hetero} & \textbf{\homo} & \textbf{\ablation} & \textbf{\hetero} & \textbf{\homo} & \textbf{\ablation} \\
        \midrule
        \relay & 1.640 & 1.777 & 1.442 & 0.684 & 0.606 & 0.592 \\
        \foraging & 2.719 & 2.913 & 1.597 & 0.784 & 0.796 & 0.558 \\
        \rendezvous & 3.019 & 3.392 & 1.682 & 0.863 & 0.902 & 0.661 \\
        \bottomrule
    \end{tabular}
    \label{tab:mi}
\end{table}

\para{MI asymmetry.} We measure the directional MI difference between slow$\to$fast and fast$\to$slow agent pairs. Across all conditions and environments, the asymmetry is near zero ($< 0.001$). There is no evidence that slow agents develop value functions that preferentially model fast agents. \textbf{H3 is not supported.}

\para{MI decay rate.} The MI decay rate (MI at lag 5 / MI at lag 0) reveals how persistent cross-agent temporal dependencies are. Heterogeneous agents show slightly higher decay rates than homogeneous agents in \relay (0.684 vs.\ 0.606), suggesting their value functions maintain cross-agent information longer across time lags. \rendezvous shows the highest decay rates overall (0.86--0.90), consistent with its strong synchronization demands. The \ablation condition consistently shows the fastest MI decay, confirming that temporal features help maintain cross-agent value function coupling.

\para{Key insight.} The \ablation condition shows 40--50\% lower MI than both \hetero and \homo conditions. This reveals that temporal context features---not decision frequency differences---are the primary driver of cross-agent value function coupling.

\subsection{Ablation: Temporal Context Features (H4)}
\label{sec:ablation}

\Tabref{tab:ablation} isolates the effect of temporal context features by comparing heterogeneous agents with and without these features.

\begin{table}[t]
    \centering
    \caption{Effect of temporal context features on heterogeneous agent performance. Features significantly help only in \rendezvous, the environment with the highest synchronization demand.}
    \begin{tabular}{@{}lcccc@{}}
        \toprule
        \textbf{Environment} & \textbf{With features} & \textbf{Without features} & \textbf{$p$-value} & \textbf{Cohen's $d$} \\
        \midrule
        \relay & $1.25 \pm 1.5$ & $1.70 \pm 0.5$ & $0.287$ & $-0.74$ \\
        \foraging & $20.67 \pm 1.2$ & $20.43 \pm 1.1$ & $0.780$ & $+0.18$ \\
        \rendezvous & $\mathbf{149.95 \pm 8.0}$ & $139.28 \pm 6.0$ & $\mathbf{0.042}$ & $\mathbf{+1.53}$ \\
        \bottomrule
    \end{tabular}
    \label{tab:ablation}
\end{table}

Temporal context features significantly improve performance only in \rendezvous ($p = 0.042$, $d = 1.53$), the environment with the highest synchronization demand. The 7.6\% improvement suggests that agents learn to use temporal awareness---normalized timestep, action frequency, and cycle phase---to better coordinate their arrival times. In \foraging, features have negligible effect ($p = 0.780$) because the task does not require synchronization. In \relay, the effect is not statistically significant ($p = 0.287$). \textbf{H4 is partially supported}: temporal context features enable coordination awareness, but only when the task demands it.

\subsection{Agent Influence Analysis}
\label{sec:influence}

We complement the MI analysis with influence scores---lagged correlations between changes in agent value functions (\eqref{eq:influence}). In \foraging, influence scores are 15--20\% higher for heterogeneous agents compared to homogeneous agents, indicating stronger inter-agent dependencies. This is consistent with the role complementarity in this environment: slow high-reward collectors and fast scouts develop value functions that are more responsive to each other's state changes. However, this stronger coupling does not manifest as \emph{hierarchical} credit assignment---the influence is symmetric rather than directional.
